\begin{table}[t]\centering\small
\caption{Examples of Prospective Technology Concepts Generated by the Proposed Framework and the
Endpoint Patents of the Predicted Links from Which They Were Derived.}
\label{tab:concept-examples}
\begin{tabularx}{\textwidth}{@{}c p{3.1cm} p{5.0cm} X@{}}
\toprule
\textbf{No.} & \textbf{Generated title} & \textbf{Endpoint patents of the predicted link} & \textbf{Generated technology concept} \\
\midrule
1 & Endoscopic imaging with virtual anatomy augmentation and 3D verification & \textbf{US-9861271-B2} \textit{System and method for tumor motion simulation and motion compensation using tracked bronchoscopy}\\[2pt]\textbf{EP-3679851-A1} \textit{Imagerie endoscopique avec parallaxe augmentée} & One method relates to a surgical imaging system comprising a manipulating device, an endoscope, and a camera. The manipulating device generates an image of the anatomical position of a patient in a surgical procedure. The imaging system is augmented to provide a virtual image of the anatomical position of a patient in a surgical procedure. In addition, the image is analyzed for accuracy by a data acquisition process. The data acquisition process may be performed remotely from the imaging system, including the acquisition of three-dimensional images of the anatomical position of objects. \\
\addlinespace[6pt]
2 & Virtual-object manipulation with sensor–camera registration in robotic surgery & \textbf{US-10383765-B2} \textit{Apparatus and method for a global coordinate system for use in robotic surgery}\\[2pt]\textbf{US-2022101745-A1} \textit{Virtual reality system for simulating a robotic surgical environment} & One or more methods of executing a surgical system are described, including the manipulation of a virtual object by a computer-aided robotic device to generate an anatomical image of a surgical scene. The system generates an image of the surgical scene based on the user's perception of the location and orientation of the surgical scene. In addition, the system provides a method for identifying and synchronizing a set of data from a sensor in the operating room with a set of data from a camera in the operating room. \\
\addlinespace[6pt]
3 & Image-based robotic planning for orthopaedic implant positioning & \textbf{US-10537388-B2} \textit{Systems and methods for planning and performing image free implant revision surgery}\\[2pt]\textbf{US-2022160440-A1} \textit{Surgical assistive robot arm} & A method of generating a robotic surgical system based on computer-aided visualization of the patient's anatomy is disclosed. The method includes techniques for adjusting the orthopaedic positioning of a surgical implant in relation to the orthopaedic features of the implant. In addition, a series of tools may be used to obtain a 3-dimensional model of a surgical procedure, such as arthroplasty, arthroplasty, and arthroplasty, to be performed within the patient's own joint space. \\
\bottomrule
\end{tabularx}
\begin{tablenotes}\footnotesize
\item Concept texts are the post-processed (grammar-corrected) outputs. Generated titles were
assigned by the authors for readability; the framework outputs concept text only.
\end{tablenotes}
\end{table}